\documentclass[a4paper]{article}
\usepackage[a4paper,margin=0mm]{geometry}
\usepackage{fontspec}
\usepackage{unicode-math}
\usepackage{tikz}
\usepackage[table]{xcolor}
\usepackage[absolute]{textpos}
\usepackage{array}
\usepackage{longtable}

\setmainfont{Libertinus Serif}
\setsansfont{Libertinus Sans}
\setmathfont{Libertinus Math}

\definecolor{accent}{HTML}{602785}
\definecolor{badge}{HTML}{B8A0C8}
\definecolor{ink}{HTML}{2A2E30}

\setlength{\TPHorizModule}{1mm}
\setlength{\TPVertModule}{1mm}
\setlength{\parindent}{0pt}
\pagestyle{empty}

\newcommand{\topicheader}[2]{%
\begin{textblock*}{194mm}(16mm,10mm)
\begin{tikzpicture}
\node[fill=accent,text=white,font=\sffamily\bfseries\fontsize{13}{13}\selectfont,
      minimum width=9mm,minimum height=9mm,inner sep=0pt] at (4.5mm,-4.5mm) {#1};
\node[font=\sffamily\bfseries\fontsize{15}{16}\selectfont,text=ink,anchor=west] at (12mm,-4.5mm) {#2};
\end{tikzpicture}
\end{textblock*}
\begin{textblock*}{194mm}(16mm,22.5mm)
\begin{tikzpicture}
\fill[accent] (0,0) rectangle (194mm,0.5mm);
\end{tikzpicture}
\end{textblock*}
}

\newcommand{\pagekicker}[1]{%
\begin{textblock*}{150mm}(16mm,25.5mm)
\begin{tikzpicture}
\node[fill=ink,text=white,font=\sffamily\bfseries\fontsize{8}{8}\selectfont,
      inner sep=2.2pt,rounded corners=1pt] {\MakeUppercase{#1}};
\end{tikzpicture}
\end{textblock*}
}

\newcommand{\pagefooter}[1]{%
\begin{textblock*}{194mm}(8mm,286.41mm)
  \begin{tikzpicture}
    \fill[accent] (0,0) rectangle (194mm,0.6mm);
  \end{tikzpicture}
\end{textblock*}
\begin{textblock*}{63.59mm}(8mm,288.27mm)
  {\sffamily\bfseries\color{accent}\fontsize{13}{16}\selectfont SAT 1200+}
\end{textblock*}
\begin{textblock*}{10mm}(188mm,288.27mm)
  \centering{\sffamily\color{accent}\fontsize{14}{16}\selectfont #1}
\end{textblock*}
}

\arrayrulecolor{accent}
\setlength{\arrayrulewidth}{0.7pt}

\begin{document}
\topicheader{23}{The Pythagorean Theorem \& Coordinates}
\pagekicker{Vocabulary for this lesson's problems (Sets 1--8)}
\pagefooter{301}
\mbox{}%
\renewcommand{\arraystretch}{1.37}
\begin{textblock*}{178mm}(16mm,33mm)
{\fontsize{11.5}{13.5}\selectfont
\begin{tabular}{|p{80mm}|p{88mm}|}
\hline
\rowcolor{accent}
{\color{white}\sffamily\bfseries English} & {\color{white}\sffamily\bfseries O'zbekcha} \\ \hline
\textbf{right triangle} & \textit{to'g'ri burchakli uchburchak} \\ \hline
\rowcolor{badge!18}
\textbf{right angle} & \textit{to'g'ri burchak} \\ \hline
\textbf{leg} & \textit{katet} \\ \hline
\rowcolor{badge!18}
\textbf{hypotenuse} & \textit{gipotenuza} \\ \hline
\textbf{shorter leg} & \textit{qisqaroq katet} \\ \hline
\rowcolor{badge!18}
\textbf{longest side} & \textit{eng uzun tomon} \\ \hline
\textbf{Pythagorean theorem} & \textit{Pifagor teoremasi} \\ \hline
\rowcolor{badge!18}
\textbf{Pythagorean triple} & \textit{Pifagor uchligi} \\ \hline
\textbf{converse (of a theorem)} & \textit{teskari tasdiq} \\ \hline
\rowcolor{badge!18}
\textbf{distance between the points} & \textit{ikki nuqta orasidagi masofa} \\ \hline
\textbf{xy-plane / origin} & \textit{xy-tekislik / koordinatalar boshi} \\ \hline
\rowcolor{badge!18}
\textbf{slope} & \textit{burchak koeffitsienti} \\ \hline
\textbf{perimeter} & \textit{perimetr} \\ \hline
\rowcolor{badge!18}
\textbf{area} & \textit{yuza} \\ \hline
\textbf{diagonal} & \textit{diagonal} \\ \hline
\rowcolor{badge!18}
\textbf{straight-line distance} & \textit{to'g'ri chiziq bo'ylab masofa} \\ \hline
\textbf{consecutive integers} & \textit{ketma-ket butun sonlar} \\ \hline
\rowcolor{badge!18}
\textbf{units apart / equally far from} & \textit{birlik uzoqlikda / teng uzoqlikda} \\ \hline
\textbf{square / rectangle} & \textit{kvadrat / to'g'ri to'rtburchak} \\ \hline
\end{tabular}
}
\end{textblock*}
\end{document}
